\section{Permutations with repetition}

Permutations without repetition assume that all objects are distinct. In many counting problems this is not true. Some objects may look the same, carry the same label, or belong to the same type. Then exchanging two identical objects does not create a new visible arrangement.

The important idea is overcounting. If we temporarily pretend that all objects are distinct, we count too many arrangements. We then divide by the number of internal swaps that cannot be seen in the final result.

\begin{definitionbox}
A permutation with repetition is an ordering of a multiset, that is, a finite collection in which some elements may occur more than once.
\end{definitionbox}

For example, the letters in the word
\[
AAB
\]
are not all distinct. The two letters \(A\) are indistinguishable. The distinct arrangements are
\[
AAB,\qquad ABA,\qquad BAA.
\]
There are three, not six.

If all three positions were occupied by distinct objects, there would be \(3!\) arrangements. But the two identical \(A\)'s can be swapped without producing a new visible arrangement. Therefore we divide by \(2!\).

\begin{theorembox}
Suppose a multiset has \(n\) elements in total, with \(n_1\) identical elements of one kind, \(n_2\) identical elements of another kind, and so on, where
\[
n_1+n_2+\cdots+n_k=n.
\]
The number of distinct permutations is
\[
\frac{n!}{n_1!n_2!\cdots n_k!}.
\]
\end{theorembox}

The numerator \(n!\) counts arrangements as if every object were distinct. The denominator removes invisible rearrangements inside each group of identical objects. If one kind appears \(n_i\) times, then those \(n_i\) identical copies can be internally rearranged in \(n_i!\) ways without changing the final visible arrangement.

\begin{examplebox}
How many distinct arrangements are there of the letters in the word \(\text{LEVEL}\)?

There are five letters. The letter \(L\) appears twice, the letter \(E\) appears twice, and the letter \(V\) appears once. Hence
\[
\frac{5!}{2!2!1!}=\frac{120}{4}=30.
\]
\end{examplebox}

\begin{examplebox}
How many distinct strings can be formed from the letters of \(\text{AABBC}\)?

There are five positions. The letter \(A\) appears twice, \(B\) appears twice, and \(C\) appears once. Therefore the number is
\[
\frac{5!}{2!2!1!}=30.
\]
\end{examplebox}

\begin{remarkbox}
Use permutations with repetition when all positions are filled, order matters, and some objects are indistinguishable. Typical examples are arrangements of letters in a word with repeated letters or arrangements of objects classified by type rather than identity.
\end{remarkbox}
